
\chapter{Карточная экосистема}\label{ch:ecosystem}
Между покупателем и~продавцом оказываются два банка, карточная сеть
и~нередко ещё процессор или PSP (payment service provider, провайдер
платёжных услуг; см.~раздел~\ref{sec:ecosystem-participants}).
Конструкция держится на межбанковской комиссии (interchange):
эмитент знает клиента и~несёт кредитный риск, эквайер подключает
продавца и~отвечает за приём, сеть передаёт сообщения, а~деньги
перемещаются не сразу и~через несколько фаз.

\section{Четырёхсторонняя модель}\label{sec:ecosystem-four-party}

Каждая карточная транзакция проходит через нескольких независимых
участников. У~каждого своя функция, свои риски и~своя экономика.

Держатель карты (cardholder) инициирует платёж. Продавец~---
торгово-сервисное предприятие, ТСП, merchant~--- принимает его.
Эмитент (issuer) отвечает за сторону плательщика, эквайер (acquirer)~---
за сторону продавца. Вместе они образуют четырёхстороннюю модель
(four-party model, иногда four-corner model), на которой построены
платёжные системы Visa, Mastercard и~Мир~\cite{visa-core-rules,mc-rules,cbr-ps}.
Над этой конструкцией стоит карточная сеть (card network, payment
scheme). Сеть задаёт правила и~поддерживает техническую инфраструктуру
для передачи сообщений и~проведения расчётов, но юридически остаётся
вне сделки между покупателем и~продавцом.

Покупатель и~продавец не могут провести карточный платёж вдвоём:
им не хватает ни масштаба, ни доверия. Продавец не проверит, есть ли
у~покупателя деньги на счёте в~другом банке; покупатель не гарантирует
оплату без участия своего банка. Между ними встают два банка: эмитент
знает покупателя и~подтверждает его платёжеспособность, а~эквайер
знает продавца и~обязуется зачислить ему деньги после расчёта.
Карточная сеть связывает их общими правилами и~технической шиной,
избавляя от необходимости договариваться напрямую. Развёрнутые
определения эмитента и~эквайера~--- в~разделе~\ref{sec:ecosystem-participants}.

\begin{figure}[tbp]
  \centering
  \includefiguresvg[width=\linewidth]{assets/figures/ch02-ecosystem/party-models}
  \caption{Четырёхсторонняя и~трёхсторонняя модели в~сравнении.}%
  \label{fig:party-models}
\end{figure}

Без карточной сети каждый из $N$ эмитентов и~$M$ эквайеров должен
был бы заключать двусторонний договор~--- $N \times M$ соглашений.
\highlight{Карточная сеть сводит эту задачу к~$N + M$}. Каждый эмитент
подписывает одно лицензионное соглашение со схемой, каждый эквайер~---
своё. С~этого момента карта любого эмитента принимается у~любого
эквайера, потому что оба подчиняются единым правилам сети. Карта Visa,
выпущенная региональным банком в~Новосибирске, работает в~терминале
кофейни в~Лиссабоне, хотя эти два банка никогда не слышали друг о~друге.

Для держателя карты и~продавца вся эта инфраструктура почти невидима.
Покупатель видит \enquote{одобрено} или \enquote{отказано}, продавец
получает деньги в~оговорённый расчётный срок. За каждым из этих
событий стоит цепочка сообщений, прошедшая от терминала через эквайера
и~карточную сеть к~эмитенту и~обратно. Каждый участник этой цепочки
принял решение, обновил свой реестр и~взял на себя часть ответственности
за результат.

\section{Участники}\label{sec:ecosystem-participants}

Поверх базовой четырёхсторонней схемы выросла прослойка посредников,
закрывающих конкретные технические или бизнес-задачи.

\subsection*{Эмитент}

Когда держатель карты прикладывает её к~терминалу, авторизационный
запрос получает эмитент и~за доли секунды решает, одобрить или
отклонить. Эмитент проверяет баланс или доступный кредитный лимит,
статус карты, географические и~временн\'{ы}е ограничения, правила
антифрода и, при необходимости, результат аутентификации~--- PIN, 3-D Secure.
Ответ \texttt{00}~--- \enquote{операция одобрена}~--- означает обязательство
оплатить эту транзакцию при клиринге; иначе продавец получает отказ,
и~покупку придётся оплатить иначе.

Экономика эмитента строится на нескольких источниках дохода: это
процентный доход по кредитным картам, годовое обслуживание, комиссия
за операции~--- снятие наличных, конвертация валюты~--- и~interchange,
межбанковская комиссия, которую эмитент получает с~каждой покупки.
Interchange определяет экономику системы и~мотивирует банки
выпускать карты, предлагать кешбэк и~мильные программы; всё это
финансируется из~комиссии, которую платит продавец. Механика
interchange разбирается в~разделе~\ref{sec:ecosystem-interchange}.

\subsection*{Эквайер}

Эмитент стоит на стороне покупателя, эквайер~--- на стороне продавца.
Эквайер подписывает с~ТСП договор эквайринга, по которому обязуется
принимать карты определённых платёжных систем, перечислять ТСП
причитающиеся суммы за вычетом комиссии и~обрабатывать чарджбэки
при оспаривании транзакции покупателем. Взамен ТСП платит
эквайеру торговую уступку (merchant service charge, MSC)~---
процент от суммы каждой транзакции, покрывающий расходы всех
участников, включая interchange и~комиссию сети.

Эквайер несёт риски, связанные с~ТСП. Если продавец обанкротится,
не доставит товар или окажется мошенником и~покупатели инициируют
чарджбэки, финансовую ответственность перед карточной сетью несёт
эквайер~--- он принял это ТСП на обслуживание. Перед подключением
эквайер проводит проверку, underwriting, и~оценивает бизнес-модель,
финансовое состояние, категорию деятельности по MCC (Merchant Category
Code, код категории ТСП по ISO~18245)~\cite{iso-18245} и~репутационные
риски. После подключения он мониторит уровень чарджбэков, подозрительные
шаблоны и~соответствие правилам схемы.

\subsection*{Карточная сеть}

Карточная сеть (или, в~терминологии 161-ФЗ,
\emph{платёжная система}) устанавливает правила, по которым работают остальные
участники, и~поддерживает техническую инфраструктуру
для передачи авторизационных сообщений, клиринга
и~расчёта~\cite{visa-core-rules,mc-rules,cbr-ps}.

Visa управляет сетью VisaNet, через которую проходят все авторизации
и~клиринговые файлы; Mastercard~--- сетью Banknet; НСПК (Национальная
система платёжных карт)~--- национальной
инфраструктурой, через которую маршрутизируются все внутрироссийские
карточные транзакции независимо от бренда на карте~\cite{nspk-about-company}.
Перевод внутрироссийского
трафика прошёл этапно в~2015~году. Mastercard переключился на ОПКЦ
(операционно-клиринговый центр НСПК)
с~1~апреля 2015~года, Visa завершила миграцию к~концу мая того же года.
Правила схемы, scheme rules, насчитывают сотни страниц и~регламентируют
детали~--- от размера логотипа на карте до порядка обработки диспутов,
от требований к~безопасности терминала до сроков расчёта между участниками.

Доход карточной сети складывается из комиссий за обработку транзакций,
оценочных сборов (assessments), лицензионных платежей участников
и~дополнительных сервисов: токенизации, скоринга антифрода, аналитики.
В~отличие от interchange, который эквайер платит эмитенту, комиссию
сети платят обе стороны: и~эмитент, и~эквайер, каждый по своей шкале.

\subsection*{Процессор}

Банки, выполняющие роли эмитента и~эквайера, далеко не всегда
обрабатывают транзакции самостоятельно. Техническую часть~---
формирование авторизационных сообщений, связь с~карточной сетью,
обработку клиринговых файлов, управление криптографическими ключами~---
часто выполняет отдельная организация, процессор (processor).

Процессоры бывают двух видов. \emph{Эмитентский процессор} (issuer
processor) обрабатывает транзакции от имени эмитента: получает
авторизационные запросы из карточной сети, сверяет их с~данными карты
и~счёта, применяет правила эмитента и~возвращает ответ. \emph{Эквайерный
процессор} (acquirer processor) работает на стороне эквайера, принимает
транзакции от терминалов, формирует сообщения в~формате карточной сети,
передаёт их на авторизацию и~обрабатывает ответы. Один банк может
использовать разных процессоров на стороне эмитента и~эквайера;
один процессор может обслуживать десятки банков.

НСПК в~эту категорию не попадает: это инфраструктурный оператор
национальной платёжной системы \enquote{Мир} и~ОПКЦ,
через который маршрутизируются все внутрироссийские карточные
транзакции~\cite{cbr-ps,nspk-about-company}. Сеть НСПК выполняет
роль карточной схемы и~центрального коммутатора; коммерческим
процессором банков она не является. До приостановки операций Visa
и~Mastercard в~России в~марте 2022~года даже покупка по карте Visa
в~магазине в~России шла через ОПКЦ НСПК (минуя VisaNet напрямую).
Независимыми процессорами
на рынке работают UCS\footnote{АО~\enquote{Компания объединённых
кредитных карточек} (КОКК), действует под брендом UCS; входит в~ТКС~Холдинг.}
(АО~КОКК, входит в~ТКС~Холдинг),
БПЦ~Процессинг (на платформе SmartVista) и~Compass~Plus
(TranzWare\slash TranzAxis).

\subsection*{Платёжный шлюз}

Платёжный шлюз (payment gateway) закрывает узкий технический фрагмент:
шифрование карточных данных, формирование сообщения для эквайера,
доставку 3-D Secure, возврат ответа на сайт продавца. Сам по себе шлюз
не подписывает с~ТСП эквайринговый договор и~не участвует в~расчётах~---
эту роль исполняет банк-эквайер, чьё подключение шлюз обслуживает.
Технические интеграторы вроде PayKeeper или Best2Pay чаще сочетают
шлюз с~функциями платёжного агрегатора и~работают по партнёрским
программам с~несколькими банками-эквайерами одновременно.

\subsection*{Провайдер платёжных услуг (PSP)}

Не каждое ТСП готово напрямую подписывать договор эквайринга с~банком,
интегрировать терминальное оборудование и~разбираться в~клиринговых
файлах. Для таких продавцов существуют провайдеры платёжных услуг
(payment service provider, PSP), берущие на себя техническую сторону
приёма платежей и~дающие продавцу единый интерфейс вместо прямой
банковской интеграции. PSP стоит между продавцом и~одним или несколькими
эквайерами, принимает платёжные данные или токен, формирует запрос
в~нужном формате, отправляет транзакцию выбранному партнёру и~возвращает
продавцу единый результат.

Объём принимаемых PSP обязательств варьируется. Когда провайдер
ограничивается технической ролью, он остаётся платёжным шлюзом.
Когда поверх шлюза появляются функции эквайера или агента, он
сам заключает договор с~ТСП и~участвует в~расчётном хвосте. На
международном рынке по этой модели работают Adyen и~Stripe;
в~России~--- ЮKassa (на стороне НКО~ЮMoney, экосистема Сбера),
CloudPayments (ООО~\enquote{КЛАУДПЭЙМЕНТС} как банковский платёжный агент Т-Банка),
Robokassa и~Best2Pay. От распределения ролей зависит, кто отвечает за
тариф, расчётные сроки, диспуты и~маршрутизацию.

\practicenote{%
  Если вы разработчик и~интегрируете приём платежей, чаще всего
  вашим контрагентом будет PSP. Вызываете его API, получаете
  асинхронное уведомление о~результате, используете его инструменты
  для возвратов и~подписок. Эквайер, карточная сеть
  и~эмитент скрыты за фасадом единого интерфейса, и~для повседневной
  работы достаточно понимать абстракцию PSP. Но когда что-то идёт
  не так (транзакция зависла, пришёл чарджбэк с~непонятным кодом
  причины, деньги не зачислились в~ожидаемый срок), знание нижележащих
  слоёв помогает найти причину и~решение.

}
\subsection*{Платёжный фасилитатор}

\textit{Платёжный фасилитатор} (payment facilitator, PayFac)~---
посредник, который подписывает с~банком-эквайером один общий договор
и~подключает под ним множество субмерчантов от своего имени, избавляя
банк от~задачи индивидуального подключения каждого ТСП.

Прямого аналога PayFac в~местном праве нет; роль реализуется через
статус банковского платёжного агента (БПА) или платёжного агрегатора
по 161-ФЗ~\cite{fz-161}. Так, АО~КОКК (UCS) выступает платёжным
агрегатором~--- БПА Т-Банка, ЮKassa агрегирует приём в~экосистеме
Сбера, а~маркетплейсы Wildberries и~Ozon строят отношения с~продавцами
по агентской модели и~замыкают приём на собственные дочерние
банки (ВБ~Банк, Озон~Банк). Полный разбор PayFac, MoR (Merchant of
Record~--- юридический продавец перед карточной сетью) и~pass-through
как разных распределений ответственности и~их соотношения с~БПА~---
в~разделе~\ref{sec:ia-payfac}.

\begin{table}[H]
  \begingroup
  \small
  \setlength{\tabcolsep}{4pt}
  \renewcommand{\arraystretch}{1.2}
  \begin{tabularx}{\textwidth}{
      @{}P{0.20\textwidth}
      Y
      Y
      Y
      Y@{}
    }
    \toprule
    \headrow
    \thd{Признак} &
    \thd{Процессор} &
    \thd{Шлюз} &
    \thd{PSP} &
    \thd{Фасилитатор} \\
    \midrule
    Лицензия & --- & --- & агрегатор & БПА \\
    Двигает деньги & нет & нет & да & да \\
    Договор с~ТСП & банк & банк & PSP & фасилитатор \\
    Риск чарджбэков & нет & нет & общий с~банком & на нём \\
    Дескриптор\footnote{Дескриптор~--- текст, который держатель карты
    видит в~выписке или мобильном банке рядом с~операцией; формируется
    стороной, чьё имя проходит в~карточной сети.} & эквайер & эквайер & ТСП & субмерчант \\
    \bottomrule
  \end{tabularx}%
  \endgroup
  \caption{Сравнение технических посредников в~карточной
  цепочке.}\label{tab:intermediaries-comparison}
\end{table}

Лицензионная строка таблицы группирует по категориям
161-ФЗ;\footnote{Собственная кредитная лицензия даёт фасилитатору
  больше свободы, чем статус БПА конкретного банка-эквайера.}
строка о~договоре~--- по тому, с~кем продавец подписывает
приём.\footnote{Договор о~приёме определяет, кто отвечает за
  тариф, сроки и~возвраты, и~часто отличается от того, чьим именем
платёж проходит в~карточной сети.} Различие между шлюзом, PSP
и~фасилитатором подвижно: один и~тот же провайдер
одновременно поставляет шлюзовую интеграцию для одних клиентов
и~работает как агрегатор для
других.\footnote{Примеры в~России (на 2026~год). \emph{Процессоры}:
  UCS (АО~КОКК), БПЦ~Процессинг, Compass~Plus; НСПК стоит вне
  категорий~--- это инфраструктурный оператор системы \enquote{Мир}
  и~ОПКЦ. \emph{Шлюзы}\slash интеграторы: PayKeeper, Best2Pay~---
  обычно совмещают с~агрегаторской ролью. \emph{PSP}: ЮKassa,
  CloudPayments, Robokassa, Best2Pay. \emph{Фасилитаторские
  конструкции}: АО~КОКК как БПА Т-Банка, агентские модели
маркетплейсов Wildberries и~Ozon.} Различия в~таблице~\ref{tab:intermediaries-comparison}
описывают модели; корпоративные ярлыки здесь вторичны.

\subsection*{Оркестратор}

Крупное ТСП работает с~несколькими эквайерами и~PSP одновременно~---
так снижается стоимость приёма, растёт доля одобрений и~обеспечивается
отказоустойчивость.

Оркестратор (payment orchestrator) управляет подключениями к~нескольким
эквайерам и~PSP и~направляет каждую транзакцию по выбранному правилу~---
по стоимости, вероятности одобрения, географии покупателя или типу карты.

От PSP и~платёжного шлюза оркестратор отличается тем, что не держит
лицензию эквайера и~не становится расчётным участником. Он
не перемещает деньги~--- он выбирает, через кого их перемещать. Регуляторная
нагрузка на оркестратор поэтому ниже, чем на PSP, но и~все финансовые
роли остаются у~подключённых через него эквайеров.

В~облачной модели (\textit{hosted orchestration}) продавец
подключается к~облачному сервису вендора, который держит коннекторы
к~эквайерам и~управляет правилами маршрутизации от своего имени.
В~локальной (\textit{self-hosted}) продавец или процессор берёт SDK
и~DSL (domain-specific language, предметно-ориентированный язык)
правил маршрутизации
и~развёртывает у~себя. Облачная требует меньше инфраструктуры на старте;
локальная даёт полный контроль над правилами и~данными транзакций.

Оркестратор окупается при работе с~несколькими активными эквайерами,
мультирегиональных выплатах или каскадных повторах с~несколькими попытками.
Главный операционный риск~--- привязка к~вендору (\textit{vendor lock-in})
на логике маршрутизации. Правила хранятся в~формате вендора, и~накопленные
годами правила трудно перенести в~другую систему без потери точности.

Оркестратор не создаёт новой банковской роли~--- он помогает управлять
уже существующими. Десять лет назад продавцу часто хватало одного
банка-эквайера, а~сегодня крупный \textit{e-commerce}-бизнес
с~международным присутствием работает с~пятью-десятью эквайерами
в~разных странах, несколькими альтернативными методами оплаты
и~множеством валют. Управлять такой схемой вручную уже невозможно.

Один участник цепочки несёт баланс и~риск, другой подписывает
договор с~ТСП, третий задаёт правила сети. Как только роли разделены,
остаётся вопрос: как между участниками делятся деньги?

\section{Кто кому платит: экономика interchange}\label{sec:ecosystem-interchange}

Эмитенты, эквайеры, карточная сеть и~процессоры стоят денег:
терминалы, обработка миллионов авторизаций в~секунду, системы антифрода,
кешбэк на карте покупателя. Распределение этих расходов между участниками
описывает экономика interchange.

\subsection*{MSC: торговая уступка}

Покупатель платит сто рублей, продавец получает не сто. Часть суммы
остаётся у~участников платежа. За право принимать карточные платежи
продавец платит эквайеру торговую уступку (merchant service charge, MSC,
иногда merchant discount rate, MDR)~--- фиксированный процент от суммы
каждой транзакции, к~которому может добавляться фиксированная плата
за операцию. MSC~--- суммарная стоимость карточного приёма для продавца.
Из ста рублей продавец получает, скажем, девяносто восемь, а~два уходят
на оплату работы всех участников цепочки.

MSC складывается из нескольких компонентов.

\begin{enumerate}
  \item \textbf{Interchange}~--- межбанковская комиссия, которую
    эквайер платит эмитенту за каждую транзакцию. Это самая
    крупная составляющая MSC.
  \item \textbf{Комиссия сети} (scheme fee, assessment)~---
    плата, которую карточная сеть взимает за обработку
    транзакции, маршрутизацию, клиринг и~расчёт. Существенно
    меньше interchange.
  \item \textbf{Маржа эквайера} (acquirer markup)~--- то, что
    остаётся эквайеру после уплаты interchange и~комиссии
    сети. Покрывает его собственные расходы на терминалы,
    процессинг, поддержку, риски, и~даёт прибыль.
\end{enumerate}

Для продавца MSC выглядит одной ставкой в~договоре, внутри
раскладывается на несколько денежных потоков. Эмитент получает свою
часть за выпуск и~обслуживание платёжного инструмента, сеть~--- за
инфраструктуру и~обработку, эквайер~--- за подключение ТСП, операционную
работу и~риск.

В~иллюстративном примере покупатель оплачивает дебетовой картой Visa
50~EUR в~супермаркете по чиповой транзакции с~MCC~5411 в~зоне действия
IFR (Interchange Fee Regulation, регламент ЕС о~межбанковских
комиссиях, 2015~год~\cite{eu-ifr}). MSC раскладывается на три
компонента~--- interchange, scheme fee, acquirer markup,~--- порядка
0,74\,\% (рис.~\ref{fig:interchange-flow}). На потребительской
кредитной карте interchange упирается в~собственный потолок IFR~---
0,30\,\%, и~MSC поднимается примерно до~0,84\,\%. Разница в~0,05~EUR
на одной транзакции выглядит ничтожной; на масштабе крупного ретейлера
с~миллионами операций она определяет, стоит ли подталкивать покупателей
к~дебетовым картам и~переключать трафик между эквайерами.

\begin{figure}[tbp]
  \centering
  \includefiguresvg[width=.95\textwidth]{assets/figures/ch02-ecosystem/interchange-flow}
  \caption{Декомпозиция MSC: дебетовая карта Visa в~IFR-зоне.}\label{fig:interchange-flow}
\end{figure}

\subsection*{Interchange: зачем эквайер платит эмитенту}

Деньги внутри схемы текут от эквайера к~эмитенту. Эмитент финансирует
выпуск карт, обслуживание счёта, антифрод, кредитный риск и~программы
лояльности; коммерческая выгода от карточного приёма возникает прежде
всего у~продавца. \highlight{Эту связку и~замыкает interchange~--- за каждую покупку
эквайер перечисляет эмитенту часть MSC}.

Направление потока~--- от эквайера к~эмитенту~---
определяется структурой издержек двустороннего рынка. Карточная сеть
работает, пока оба конца цепочки активны~--- покупатели носят карты,
продавцы их принимают. Продавец готов платить за карточный приём,
потому что карта увеличивает средний чек и~устраняет расходы
на инкассацию. У~эмитента~--- свой набор расходов: выпуск и~перевыпуск
пластика, процессинг авторизаций, риск невозврата по кредитным картам,
обслуживание споров и~--- если банк хочет конкурировать за держателя~---
кешбэк и~мильные программы. Без interchange единственные источники дохода эмитента
по дебетовой карте~--- годовое обслуживание и~комиссия за снятие
наличных; этого не хватает, чтобы окупить инфраструктуру массовой
эмиссии. Банки выпускают меньше карт, покупатели реже платят картой,
продавцы видят меньше карточных транзакций, сеть сужается. Interchange
разрывает этот цикл, перенося часть выручки от стороны, получающей
коммерческую выгоду через продавца и~эквайера, к~стороне с~наибольшими
инфраструктурными расходами~--- эмитенту.

Альтернативный дизайн~--- прямые платежи от продавца эмитенту без
посредника~--- потребовал бы $N \times M$ двусторонних соглашений, ровно
той комбинаторной проблемы, которую карточная сеть решает для маршрутизации.
Interchange работает по тому же принципу: сеть задаёт единую таблицу
ставок, и~каждый участник подчиняется ей без переговоров с~каждым
контрагентом.

Когда Европейский союз ограничил interchange Регламентом IFR
уровнем 0,2\,\% для дебетовых и~0,3\,\% для кредитных карт,
последствия проявились на обоих концах цепочки~\cite{eu-ifr}.
Эмитенты урезали программы лояльности,
перешли на годовые комиссии за обслуживание и~сосредоточили кешбэк
на премиальных продуктах, не подпадающих под регулирование,~--- коммерческих
картах и~трёхсторонних схемах. Продавцы получили снижение MSC, но
не всегда пропорциональное~--- часть экономии осталась у~эквайеров в~виде
возросшей маржи.

Размер interchange задаёт карточная сеть; свободные переговоры
между конкретным эмитентом и~эквайером тут не работают. Ставка
зависит от типа карты: дебетовая обычно дешевле кредитной.
Зависит от категории ТСП по MCC: продуктовый магазин платит меньше,
чем ювелирный. Зависит от способа приёма: контактная транзакция
дешевле ручного ввода номера карты. Зависит от географии: внутренняя
транзакция дешевле международной. Полные таблицы interchange
публикуются на сайтах Visa и~Mastercard~--- сотни строк с~различными
комбинациями~\cite{visa-interchange,mc-interchange-us}.

\importantnote{%
  Продавец платит эквайеру MSC, а interchange живёт внутри этой суммы~---
  это часть MSC, которую
  эквайер перечисляет эмитенту. Продавец может даже не знать,
  какую долю его MSC составляет interchange, если договор
  с~эквайером предусматривает единую ставку, flat rate,
  вместо модели interchange++, при которой продавец видит interchange
  и~маржу эквайера отдельно. Когда продавец \enquote{переплачивает
  за эквайринг}, проблема часто в~высоком interchange, на который
  эквайер повлиять не может.

}
\subsection*{Регулирование interchange}

Interchange остаётся предметом постоянных споров между торговыми
сетями, банками и~регуляторами. Продавцы считают ставки завышенными;
банки отвечают, что резкое снижение interchange подорвёт программы
лояльности и~сократит стимулы к~эмиссии карт; регуляторы пытаются
найти баланс.

В~Европейском союзе спор привёл к~уже упомянутому Регламенту IFR (2015),
ограничившему interchange
для потребительских дебетовых карт уровнем 0,2\,\% от суммы транзакции,
для кредитных~--- 0,3\,\%~\cite{eu-ifr}. Регламент задаёт потолок,
но сам по себе не определяет, как рынок перераспределит последствия
между продавцом, эквайером, эмитентом и~держателем карты.

В~161-ФЗ такого общего потолка для interchange
не установлено~\cite{fz-161}.

\subsection*{Модель ценообразования: что видит продавец}

Для продавца interchange важен и~как абстрактный механизм, и~как часть
тарифной модели. В~договоре эквайринга его показывают одним из трёх способов.

\begin{itemize}
  \item \textbf{Flat rate}~--- единая ставка MSC на все
    транзакции, например 2,0\,\%. Просто, предсказуемо,
    но невыгодно для продавцов с~высокой долей дебетовых карт,
    по которым interchange ниже.
  \item \textbf{Interchange++}~--- продавец видит фактический
    interchange по каждой транзакции, комиссию сети и~наценку
    эквайера отдельно. Прозрачно, даёт возможность оптимизировать,
    но требует экспертизы для анализа выписок.
  \item \textbf{Tiered pricing}~--- эквайер группирует транзакции
    в~категории \enquote{квалифицированные},
    \enquote{среднеквалифицированные} и~\enquote{неквалифицированные}
    с~разными ставками. Непрозрачно и~постепенно уходит с~рынка.
\end{itemize}

Выбор модели определяет, насколько легко продавцу понять причину своих
издержек, сопоставить тариф с~качеством одобрения и~принять решение
о~маршрутизации трафика между провайдерами.

\practicenote{%
  Если вы проектируете платёжную аналитику для ТСП или платформу
  управления платежами, поддержка модели interchange++ с~детализацией
  до уровня отдельной транзакции нужна большинству крупных продавцов.
}
\section{Трёхсторонняя модель}\label{sec:ecosystem-three-party}

Альтернатива четырёхсторонней схеме~--- трёхсторонняя модель,
three-party model, closed-loop model, в~которой одна организация
совмещает роли эмитента, эквайера и~карточной сети. Классический
пример~--- American Express, Amex.

В~чистой трёхсторонней модели Amex сам выпускает карты, сам подписывает
договоры с~продавцами и~сам обрабатывает транзакции от начала до конца.
Interchange здесь теряет смысл~--- нет двух банков, между которыми нужно
передавать комиссию, и~вся выручка от продавца, merchant discount,
остаётся внутри одной организации. Amex получает полный контроль над
клиентским опытом на обеих сторонах, и~для держателя карты, и~для
продавца. Чтобы карта работала в~новом магазине, Amex должен
подписать договор с~ним напрямую. В~четырёхсторонней модели достаточно,
чтобы у~магазина был любой эквайер, подключённый к~сети. Структурное
сравнение двух моделей~--- на~рис.~\ref{fig:party-models}.

Ограничение масштаба объясняет, почему четырёхсторонняя модель
доминирует. Visa и~Mastercard присутствуют практически в~каждой
торговой точке мира, тогда как Amex исторически сосредоточен
в~премиальном сегменте с~более ограниченным покрытием. Публичные
материалы American Express прямо описывают сеть как инфраструктуру
для самой компании, сторонних эмитентов и~эквайеров~--- классический
closed-loop давно дополняется партнёрской
сетью~\cite{amex-network}.

Четырёхсторонняя схема лучше масштабирует приём и~специализацию ролей;
трёхсторонняя даёт более плотное управление всем циклом операции.
